\chapter{迴圈}
\label{ch:loops}
\index{迴圈}

生活裡有兩種「重複做同一件事」：

\begin{itemize}
\item 第一種你一開始就知道要做幾次，摺完一籃衣服，你知道今天總共要摺幾件；打掃前，你
知道家裡有幾扇窗戶要擦。
\item 第二種你不知道要做幾次，只知道「做到某個狀態出現為止」，洗碗要洗到碗槽空了才能
停，存錢要存到存夠目標金額才能停，你沒辦法在動手之前就先算出精確的次數。
\end{itemize}

這兩種重複，光靠你之前學過的工具都解決不了，你不可能把\texttt{print} 打五十次來擦五十
扇窗戶，更不可能在動手存錢之前，先算好到底要存幾天。這一章要教的，就是 Python 怎麼分別
處理這兩種重複，讓你只寫一次程式碼，就能重複執行任意多次。

這一章分成幾段推進：先教「已知次數」的 \texttt{for}，再教「不知道要做幾次、要看狀態
才知道該不該繼續」的 \texttt{while}。中間會穿插兩個實用的補充技巧：怎麼在一個迴圈裡面
再放一個迴圈，處理「每一組裡的每一個」這種雙層重複；以及怎麼在走訪一串東西的同時，順便
拿到「這是第幾個」的位置資訊。

接著回頭處理幾件更細節的事：怎麼在迴圈進行到一半時喊停或
跳過某一輪，以及怎麼知道一個迴圈到底是正常跑完、還是被提早喊停結束的。最後，把前面所有
例子拉高一層來看，你會發現不管是擦窗戶、存錢還是翻抽屜找鑰匙，背後其實只有三種常見的
迴圈任務「長相」：累加、計數、搜尋。

這三種長相會在你之後寫的幾乎每一支程式裡反覆出現，
提早認出它們的骨架，之後遇到新的問題時，會更容易想到「這其實就是搜尋的模式」，不用每次
都從零開始想解法。這一章跟前面幾章一樣是地基章節，技巧點不少，但每一個都值得你自己動手
重現一次，光是看得懂並不算數。

\section{已知次數的重複：\texttt{for} 搭配 \texttt{range()}}
\label{sec:for-range}

先從最單純的情境開始：你動手之前就知道要做幾次，不涉及任何「這次要不要繼續」的判斷，
只是把同一個動作重複套用固定的次數。Python 用 \texttt{for} 搭配一個叫 \texttt{range()}
的工具處理這種情境。\texttt{range(n)} 會產生一串數字，讓 \texttt{for} 迴圈依序拿去用，
剛好重複 \texttt{n} 次。
\index{for（搭配 range()）}

\begin{code}{python}{for 搭配 range() 重複已知次數}
for i in range(n):
    repeated_action
\end{code}

\texttt{range(n)}裡的\texttt{n}就是你想重複的次數；\texttt{i}是每一輪迴圈都會拿到的
一個數字，你可以在迴圈裡用它，也可以完全不用它。重點是縮排在\texttt{for}底下的那幾行，
會被原封不動地重複執行\texttt{n}次。

\begin{example}{python}{(windows.py) 擦五扇窗戶}
for i in range(5):
    print(f"Cleaning window {i + 1}")

print("All windows cleaned!")
\end{example}

\iosample{}{Cleaning window 1 \\ Cleaning window 2 \\ Cleaning window 3 \\ Cleaning window 4
\\ Cleaning window 5 \\ All windows cleaned! \\}

留意這裡多寫了一個\texttt{+ 1}。這不是隨便加的，是為了修正一個新手幾乎都會踩到的坑。
把上面例子裡的\texttt{i + 1}換回單純的\texttt{i}，實際印出來的會是這樣：

\begin{example}{python}{(range-check.py) 直接印出 i}
for i in range(5):
    print(i)
\end{example}

\iosample{}{0 \\ 1 \\ 2 \\ 3 \\ 4 \\}

\keypoint{\texttt{range(5)}確實會讓迴圈跑滿 5 次，但拿到的數字是從 0 開始、到 4 為止，
不是從 1 到 5。很多人第一次看到會誤以為最後應該印出 5，其實 5 從來沒有出現過。這是
\texttt{range()}最容易讓人誤會的地方：次數是對的，起點卻不是你直覺以為的 1。想要顯示
「第幾次」給人看的時候（像上面擦窗戶的例子），記得自己補一個\texttt{+ 1}。}

除了只給一個數字，\texttt{range()}還可以指定起點、終點、跟每次要跳幾格（間隔），間隔
也可以是負的，讓數字倒著數。

\begin{code}{python}{range() 指定起點、終點與間隔}
for i in range(start, stop, step):
    repeated_action
\end{code}

如果只給兩個數字\texttt{range(start, stop)}，Python 會假設間隔是 1；如果間隔是負的，
數字就會從\texttt{start}開始往下數，一路數到比\texttt{stop}大一點點的地方為止。

\begin{example}{python}{(squats.py) 深蹲倒數}
for remaining in range(10, 0, -1):
    print(f"Reps remaining: {remaining}")

print("Squats complete!")
\end{example}

\iosample{}{Reps remaining: 10 \\ Reps remaining: 9 \\ Reps remaining: 8 \\ Reps remaining: 7
\\ Reps remaining: 6 \\ Reps remaining: 5 \\ Reps remaining: 4 \\ Reps remaining: 3 \\
Reps remaining: 2 \\ Reps remaining: 1 \\ Squats complete! \\}

這裡\texttt{range(10, 0, -1)}是「從 10 開始，每次減 1，數到比 0 大的地方為止」，所以
最後印出來的是 1，不是 0。跟前面\texttt{range(n)}的道理一樣，終點那個數字本身永遠不會
被包含進去。

\begin{drills}[label=drills:for-range]
\item 不要真的執行，用眼睛看：\texttt{for i in range(3): print(i)}會印出哪三行？
\item \texttt{for i in range(2, 9, 2): print(i)}最後一行印出的數字是多少？
\item 填空：想要一個迴圈從 9 開始、每次減 2、數到 1（含）為止，\texttt{range()}該怎麼
填？（提示：終點那個數字不會被包含進去，所以終點要填比 1 還小的數字。）
\item \texttt{for i in range(5, 5): print(i)}這段程式碼，迴圈本體總共會執行幾次？
\end{drills}

\section{對一串東西的每一個都做一次：\texttt{for} 走訪字串}
\label{sec:for-string}

前一節的\texttt{for}是「把同一個動作重複做\texttt{n}次」，這一節要教的\texttt{for}
語感不太一樣，重點換成「對一串東西裡的每一個都做一次」。這串東西可以是你在第 \ref{ch:strings} 章學過的
字串。字串本身就是一串字元排在一起，\texttt{for}可以直接一個一個走訪它們，不用先算出
字串有幾個字元、也不用自己管索引數字。
\index{for 走訪序列（例如字串）}

\begin{code}{python}{for 走訪字串的每個字元}
for char in text:
    repeated_action
\end{code}

每一輪迴圈，\texttt{char}會依序變成\texttt{text}裡的下一個字元，直到走完整個字串為止。

\begin{example}{python}{(count-spaces.py) 數一句話裡有幾個空白字元}
sentence = "nice weather today"
space_count = 0

for char in sentence:
    if char == " ":
        space_count += 1

print(f"This sentence has {space_count} spaces")
\end{example}

\iosample{}{This sentence has 2 spaces \\}

留意這裡的\texttt{space\_count += 1}，這是\texttt{space\_count = space\_count + 1}
的簡寫，意思是「把\texttt{space\_count}加 1 的結果，重新貼回\texttt{space\_count}這個
名字牌」，呼應第 \ref{ch:variables} 章教過的名字牌模型。之後你會常看到的\texttt{-=}則是減法版本，效果
一樣是簡寫，不是新的運算方式。
\index{複合賦值運算子（+=／-=）}

\begin{example}{python}{(check-letter.py) 檢查一個字裡有沒有特定字母}
word = "python"
target = "y"
found = False

for letter in word:
    if letter == target:
        found = True

if found:
    print(f"The word '{word}' contains the letter '{target}'")
else:
    print(f"The word '{word}' does not contain the letter '{target}'")
\end{example}

\iosample{}{The word 'python' contains the letter 'y' \\}

\stylepoint{如果你只是想拿到字元本身，直接\texttt{for char in word}走訪就好，不需要先
繞一圈去算索引數字（例如\texttt{for i in range(len(word))}再用\texttt{word[i]}取出
字元）。除非你真的需要知道「這是第幾個字元」這個位置資訊，不然直接走訪內容本身才是更
Python 的寫法；如果真的需要位置資訊，本章稍後會介紹更好的做法（\texttt{enumerate()}），
不用回頭走\texttt{range(len(word))}這條路。}

之後你學到 list，會發現這個\texttt{for}寫法幾乎不用改，list 走訪的方式跟字串一模一樣，
只是走訪的對象換成另一種容器。

\begin{drills}[label=drills:for-string]
\item 不要真的執行，用眼睛看：\texttt{for char in "cat": print(char.upper())}會印出
哪三行？
\item 抓蟲：下面這段程式碼執行後會出現錯誤，錯誤訊息會是什麼？該怎麼修？
\begin{verbatim}
for char in "abc"
    print(char)
\end{verbatim}
\end{drills}

\section{迴圈裡面還有迴圈：巢狀迴圈}
\label{sec:nested-loops}

前面兩節的\texttt{for}走訪的都是「一層」的重複，不管是固定次數還是走訪字串，每一輪
迴圈做的都是一件單純的事。但生活裡的重複性任務常常不只一層。

打掃房子的時候，你面對的
其實是「對每一個房間，重複擦每一扇窗戶」這種雙層任務：外層是一個一個房間輪流處理，內層
是處理每個房間時，一扇一扇窗戶擦過去。這種「迴圈裡面還有一個迴圈」的寫法，叫做巢狀迴圈，
寫法就是把一個\texttt{for}迴圈整段，直接放進另一個\texttt{for}迴圈的縮排底下。
\index{巢狀迴圈}

\begin{code}{python}{巢狀迴圈：for 裡面再放一個 for}
for outer_item in outer_sequence:
    repeated_action
    for inner_item in inner_sequence:
        repeated_action
\end{code}

外層迴圈每跑一輪，內層迴圈就會完整跑過一次自己全部的輪次；外層跑幾輪，內層那一整套動作
就會被重複幾次。縮排在這裡特別重要：內層\texttt{for}縮排比外層多一層，Python 才認得出
「這是外層迴圈裡面的東西」，不是外層迴圈外面的獨立程式碼。

\begin{example}{python}{(clean-rooms.py) 打掃每個房間的每一扇窗戶}
for room in range(1, 4):
    print(f"Entering room {room}")
    for window in range(1, 3):
        print(f"  Cleaning room {room}, window {window}")
\end{example}

\iosample{}{Entering room 1 \\ ~~Cleaning room 1, window 1 \\ ~~Cleaning room 1, window 2
\\ Entering room 2 \\ ~~Cleaning room 2, window 1 \\ ~~Cleaning room 2, window 2
\\ Entering room 3 \\ ~~Cleaning room 3, window 1 \\ ~~Cleaning room 3, window 2 \\}

外層\texttt{room}先變成 1，進入第 1 個房間之後，內層\texttt{window}迴圈完整跑了 1、2
兩輪，把第 1 個房間的兩扇窗戶都擦完，才會回到外層，讓\texttt{room}變成 2，重新進入內層
迴圈再跑一次。整個迴圈總共會執行 3 乘 2 等於 6 次擦窗戶的動作，不是 3 加 2 次。

巢狀迴圈有一個新手很容易誤會的地方：\texttt{break}或\texttt{continue}只會作用在離它
最近的那一層迴圈上，不會連帶影響外層迴圈。

\begin{example}{python}{(clean-rooms-break.py) break 只跳出內層迴圈}
for room in range(1, 3):
    print(f"Cleaning room {room}")
    for window in range(1, 4):
        if window == 2:
            break
        print(f"  Finished window {window}")
    print(f"Room {room} cleaning done")
\end{example}

\iosample{}{Cleaning room 1 \\ ~~Finished window 1 \\ Room 1 cleaning done \\ Cleaning room 2
\\ ~~Finished window 1 \\ Room 2 cleaning done \\}

留意每個房間都只印出「Finished window 1」，第 2、3 扇窗戶完全沒被擦到，因為\texttt{window}
一等於 2，\texttt{break}就讓內層\texttt{for}迴圈提早結束了。但外層\texttt{for}迴圈完全
沒受影響，兩個房間都照常輪到，各自印出「cleaning done」。如果你原本期待\texttt{break}能讓整個
巢狀迴圈（包含外層）都停下來，這裡就是一個容易踩到的坑。

\keypoint{\texttt{break}／\texttt{continue}的作用範圍，只到離它最近的那一層迴圈為止。
想讓外層迴圈也跟著提早結束，得在外層迴圈自己的層級再加一次判斷（例如另外設一個變數記錄
「要不要繼續跑外層」），單靠內層的\texttt{break}辦不到。}

巢狀迴圈也可以搭配第~\ref{sec:accumulate-count}~節會正式介紹的累加骨架一起用：外層跟
內層各跑各的輪次，但用同一個變數，把每一輪的結果一路加總起來。

\begin{example}{python}{(total-windows.py) 算出整棟房子總共有幾扇窗戶}
total_windows = 0

for room in range(1, 4):
    for window in range(1, 3):
        total_windows += 1

print(f"The whole house has {total_windows} windows total")
\end{example}

\iosample{}{The whole house has 6 windows total \\}

這裡\texttt{total\_windows}在兩層迴圈外面先初始化成 0，不管走到外層第幾個房間、內層第
幾扇窗戶，每一輪都對同一個\texttt{total\_windows}做\texttt{+= 1}，兩層迴圈跑完之後，
\texttt{total\_windows}就會是 3 個房間乘 2 扇窗戶，總共 6 扇窗戶的總和，不用另外拆開來
一個房間一個房間手動加。

\begin{drills}[label=drills:nested-loops]
\item 不要真的執行，用眼睛看：下面這段程式碼總共會印出幾行「check」？
\begin{verbatim}
for i in range(2):
    for j in range(3):
        print("check")
\end{verbatim}
\item 承接上面\texttt{clean-rooms.py}：如果外層\texttt{range(1, 4)}改成
\texttt{range(1, 6)}，內層\texttt{range(1, 3)}不變，整個迴圈總共會印出幾次「Cleaning」
開頭的訊息？
\item 抓蟲：如果你想讓\texttt{clean-rooms-break.py}裡的\texttt{break}讓兩層迴圈都一起
結束，只在內層寫\texttt{break}是不夠的，你還需要什麼？（不用寫出完整程式碼，說明思路
即可）
\end{drills}

\section{同時要位置又要內容：\texttt{enumerate()} 搭配 \texttt{for}}
\label{sec:enumerate}

第~\ref{sec:for-string}~節教過，\texttt{for char in word}這種寫法能直接拿到字串裡的
每個字元本身，不用自己管索引數字。但有些情況你兩個都要：不只要知道「這個字元是什麼」，
還要知道「它排在第幾個位置」。例如你想知道某個字母第一次出現在字串的哪個位置，只走訪
字元本身是不夠的，你還得同時追蹤位置。Python 提供一個現成的工具處理這種「內容跟位置都
要」的情境，叫\texttt{enumerate()}。
\index{enumerate()}

\begin{code}{python}{for 搭配 enumerate() 同時取得位置與內容}
for index, item in enumerate(sequence):
    repeated_action
\end{code}

\texttt{enumerate(sequence)}把\texttt{sequence}包成一種每一輪都會吐出「位置、內容」
這一組配對的東西，\texttt{for}迴圈這裡一次接住兩個名字：\texttt{index}是這一輪走到第幾
個位置（從 0 開始算，跟\texttt{range()}的起點一樣），\texttt{item}則是這個位置上的內容
本身。之後在講 tuple 的那一章，你會學到這種一次接住好幾個值的寫法背後的完整機制，這裡
先掌握「\texttt{enumerate()}同時給你位置跟內容」這個用法就夠了。

\begin{example}{python}{(find-letter.py) 找出某個字母第一次出現的位置}
word = "banana"
target = "n"

for index, char in enumerate(word):
    if char == target:
        print(f"'{target}' first appears at index {index}")
        break
\end{example}

\iosample{}{'n' first appears at index 2 \\}

迴圈依序走到索引 0（\texttt{b}）、1（\texttt{a}）都不符合，走到索引 2（\texttt{n}）
才符合\texttt{target}，印出訊息後就用\texttt{break}提早結束，後面的索引 3、4、5 都不會
再檢查。留意這裡把\texttt{enumerate()}跟第~\ref{sec:break-search}~節的線性搜尋合在
一起用：一邊走訪、一邊記錄位置，一找到符合的字元就用\texttt{break}提早結束，是很常見的
組合。

\begin{example}{python}{(enumerate-oops.py) 忘記同時接住兩個值}
word = "python"

for index in enumerate(word):
    print(index)
\end{example}

\iosample{}{(0, 'p') \\ (1, 'y') \\ (2, 't') \\ (3, 'h') \\ (4, 'o') \\ (5, 'n') \\}

如果只寫一個名字\texttt{index}去接\texttt{enumerate(word)}吐出來的東西，Python 不會
報錯，但\texttt{index}接到的其實是「位置跟內容綁在一起」的一整組配對，印出來會是像
\texttt{(0, 'p')}這樣一組一組的東西，跟原本想要的乾淨數字或字元不一樣。想要兩個分開拿，
就得像前面\texttt{find-letter.py}那樣，一次寫兩個名字\texttt{index}、\texttt{char}
去接。

\stylepoint{只要你需要「這是第幾個」這個位置資訊，優先用\texttt{enumerate()}，不要自己
另外造一個變數手動累加計數。手動維護一個計數器正是下一節要比較的\texttt{while}常見的
負擔，\texttt{enumerate()}讓\texttt{for}也能拿到位置資訊，卻不用揹上這個負擔。}

\begin{drills}[label=drills:enumerate]
\item 不要真的執行，用眼睛看：下面這段程式碼會印出哪幾行？
\begin{verbatim}
for index, char in enumerate("cat"):
    print(f"{index}: {char}")
\end{verbatim}
\item 填空：想要印出\texttt{word = "hello"}裡每個字母的「第幾個」（從 1 開始算，不是
從 0），迴圈該怎麼寫？（提示：跟第~\ref{sec:for-range}~節的\texttt{+ 1}是同樣的手法。）
\item 抓蟲：下面這段程式碼想找出\texttt{"apple"}裡字母\texttt{"p"}第一次出現的位置，
但實際執行會印出兩個索引，問題出在哪裡？
\begin{verbatim}
word = "apple"

for index, char in enumerate(word):
    if char == "p":
        print(index)
\end{verbatim}
（提示：對照\texttt{find-letter.py}，找到之後少了什麼？）
\end{drills}

\section{不確定次數的重複：\texttt{while}}
\label{sec:while-basics}

前面幾節的\texttt{for}，不管是重複固定次數、走訪一串東西，還是巢狀著用、搭配
\texttt{enumerate()}一起拿位置資訊，都有一個共同點：Python 在迴圈開始之前，其實就已經
知道要跑幾輪。

但生活裡更常見的重複其實不是這樣，洗碗要洗到碗槽空了才能停，存錢要存到
存夠目標金額才能停，你沒辦法在動手之前就先算出精確的次數，只知道「什麼狀態出現時該停
下來」。這種「不確定要做幾次，但知道停止條件」的重複，要用\texttt{while}處理。
\index{while}

\begin{code}{python}{while 迴圈：先檢查、再決定要不要做}
while condition:
    repeated_action
    update_variable_used_in_condition
\end{code}

\texttt{while}每一輪開始前都會先檢查\texttt{condition}，條件成立才執行縮排底下的內容，
執行完再回頭檢查一次，條件不成立了才停下來離開迴圈。

\begin{example}{python}{(saving.py) 存錢存到存夠目標金額}
savings = 0
daily_amount = 45
target = 200
days = 0

while savings < target:
    savings += daily_amount
    days += 1

print(f"Saved for {days} days, total saved: {savings} dollars")
\end{example}

\iosample{}{Saved for 5 days, total saved: 225 dollars \\}

留意最後存到的金額是 225 元，比目標的 200 元還多。這很正常，跟現實生活一樣：
\texttt{while}只在乎「這一輪開始前，條件成不成立」，不會幫你精準停在剛好等於目標的那
一刻。存到第 5 天時，存款從 155 元加到 200 元，已經到達目標了，但這一輪動作已經做完，
程式才會在下一輪開始前重新檢查，發現\texttt{200 < 200}不成立，於是停下來，不是在存錢的
當下就精準喊停。

\texttt{while}迴圈最容易出的問題，是忘記更新\texttt{condition}用到的變數，導致條件
永遠成立、迴圈永遠不會停。
\index{無窮迴圈}

\begin{example}{python}{(saving-bug.py) 忘記更新條件變數}
savings = 0
target = 100

while savings < target:
    print("Saved a little money")
\end{example}

這段程式碼裡，\texttt{savings}從頭到尾都是 0，\texttt{savings < target}永遠成立，
迴圈永遠不會停下來。實際執行的話，螢幕會被「Saved a little money」這行字灌爆，程式看起來像當機了
一樣，其實只是陷入了無窮迴圈。想要停下來，只能按下\texttt{Ctrl+C}（或是在編輯器裡按下
停止鍵）強制中斷程式。修法是在迴圈裡把\texttt{savings}真的加上去，例如補上
\texttt{savings += 50}這一行，讓條件變數每一輪都真的往目標靠近。

\keypoint{\texttt{while}迴圈能不能停下來，完全取決於迴圈裡有沒有確實更新
\texttt{condition}用到的變數。寫\texttt{while}迴圈時，養成習慣問自己一句：「這個條件
會不會在某一輪之後變成不成立？」如果答案是不會，這個迴圈就永遠不會停。}

忘記更新變數只是無窮迴圈最常見的一種寫法，還有一種更隱蔽的版本：變數確實有更新，但更新
的方式讓條件永遠不會變成不成立。

\begin{example}{python}{(never-zero.py) 條件用 != 卻永遠等不到那個值}
n = 5

while n != 0:
    print(n)
    n -= 2
\end{example}

\texttt{n}從 5 開始，每一輪減 2：5、3、1，再減 2 變成 -1、-3、-5……\texttt{n}永遠不會
剛好等於 0，因為它是用「跳 2 格」的方式往下數，直接跳過了 0 這個值。條件\texttt{n != 0}
（\texttt{n}不等於 0）就永遠成立，即使變數確實每一輪都有更新，這仍然是一個貨真價實的
無窮迴圈。

\keypoint{拿\texttt{!=}當\texttt{while}的停止條件時要特別小心：只要變數更新的方式
（例如每次加 2、減 3）有可能整個跳過你比較的那個值，條件就永遠不會變成不成立。想確保
迴圈一定停得下來，優先用\texttt{<}或\texttt{>}這種「範圍」型的比較，不要用\texttt{!=}
這種「剛好等於」型的比較。像\texttt{savings < target}或\texttt{n > 0}，就算變數不小心
跳過某個精確數字，範圍式的條件還是會在下一輪就抓到「已經不成立」。}

\begin{drills}[label=drills:while-basics]
\item 不要真的執行，用眼睛看：下面這段程式碼會印出幾次「chances left」？
\begin{verbatim}
n = 3
while n > 0:
    print(f"{n} chances left")
    n -= 1
\end{verbatim}
\item 抓蟲：下面這段程式碼會陷入無窮迴圈，問題出在哪？該怎麼改？
\begin{verbatim}
attempts = 3

while attempts > 0:
    print(f"{attempts} chances left")
\end{verbatim}
\item 抓蟲：下面這段程式碼也會陷入無窮迴圈，問題出在哪？（提示：注意迴圈裡到底更新的是
哪一個變數）
\begin{verbatim}
remaining = 3
count = 0

while remaining > 0:
    print(f"{remaining} chances left")
    count += 1
\end{verbatim}
\item 承接\texttt{never-zero.py}：如果把\texttt{n = 5}改成\texttt{n = 6}，其他都不變，
這個迴圈還會是無窮迴圈嗎？為什麼？
\end{drills}

\section{不知道要問幾次：\texttt{while} 搭配 \texttt{input()}}
\label{sec:while-input}

前一節的存錢範例，\texttt{while}要不要繼續是靠程式自己算出來的；但很多時候，決定迴圈
要跑幾輪的人其實是使用者，你不知道他今天要記錄幾件家事，只知道他打完之後，會告訴你
「好了」。還記得第 \ref{ch:strings} 章用過的\texttt{input()}嗎？這裡我們要反覆用它讀取使用者輸入，
每一輪迴圈都問一次，直到使用者輸入約定好的結束關鍵字為止。

\begin{code}{python}{while 搭配 input() 反覆詢問}
value = input(prompt_text)

while value != stop_keyword:
    repeated_action
    value = input(prompt_text)
\end{code}

留意這裡\texttt{input()}被呼叫了兩次：一次在迴圈開始之前，先問一輪、拿到第一個答案；
一次在迴圈裡面，每做完一輪就再問一次、拿新的答案回來，讓下一輪開始前的條件檢查看到的是
最新的輸入。

\begin{example}{python}{(chores.py) 反覆記錄今天做了哪些家事}
chore = input("What chore did you do today? (Type 'done' to finish)")

while chore != "done":
    print(f"Logged: {chore}")
    chore = input("Anything else? (Type 'done' to finish)")

print("Chore log complete!")
\end{example}

\iosample{wash dishes \\ take out trash \\ done \\}{Logged: wash dishes \\ Logged: take out trash \\ Chore log complete! \\}

這裡只用到\texttt{input()}最基本的能力，拿到使用者打的一行字，拿去跟一個固定的字串
比較。至於提示文字要怎麼設計得更友善、或是把使用者輸入的內容轉成數字來用，這本書之後
需要的時候會再回頭補充，這裡先掌握「\texttt{while}的停止條件可以來自使用者輸入」這件
事就夠了。

\begin{drills}[label=drills:while-input]
\item 承接上面的\texttt{chores.py}：如果使用者依序輸入「mop floor」、「done」，程式總共會
印出幾行「Logged:」？
\item 如果使用者一開始就直接輸入「done」，程式還會印出任何一行「Logged:」嗎？
\end{drills}

\section{\texttt{for} 還是 \texttt{while}？怎麼選}
\label{sec:for-vs-while}

現在你已經看過\texttt{for}的幾種用法（固定次數、走訪字串、巢狀、搭配
\texttt{enumerate()}），也看過\texttt{while}怎麼處理不確定次數的重複。兩種語法都能寫出
「重複做某件事」的程式，新手常見的疑惑是：這兩個工具，到底該挑哪一個？

判斷的關鍵只有一個問題：動手之前，你知不知道要重複幾次，或者知不知道要處理的是哪一批
東西？知道的話，用\texttt{for}；只知道「什麼狀態出現時該停」，但不知道要跑幾輪，才用
\texttt{while}。擦五扇窗戶、走訪一句話裡的每個字元，這些你在動手之前就已經確定範圍，
用\texttt{for}最直接。存錢存到存夠、反覆詢問使用者直到他說「好了」，這些你不知道要跑
幾輪，只能用\texttt{while}邊做邊檢查。

實務上可以用兩個問題快速判斷：

\begin{itemize}
\item 動手之前，「要跑幾輪」或「要處理哪些東西」是不是已經確定？確定的話，用
\texttt{for}。
\item 停止的判斷，是不是要看執行過程中才會出現、事先算不出來的狀態（使用者說了什麼、
存款有沒有到達目標）？是的話，用\texttt{while}。
\end{itemize}

這兩個問題的答案通常會一致：已知範圍幾乎都對應「用\texttt{for}」，看狀態才能決定幾乎都
對應「用\texttt{while}」。真正會讓人猶豫的情境並不多，多半是因為把一個其實已知的範圍，
誤看成了不確定的狀態，或是反過來。

\begin{example}{python}{(windows-while.py) 把擦窗戶的例子改寫成 while}
window = 0

while window < 5:
    print(f"Cleaning window {window + 1}")
    window += 1

print("All windows cleaned!")
\end{example}

\iosample{}{Cleaning window 1 \\ Cleaning window 2 \\ Cleaning window 3 \\ Cleaning window 4
\\ Cleaning window 5 \\ All windows cleaned! \\}

這段程式碼跟第~\ref{sec:for-range}~節的\texttt{windows.py}做的事完全一樣，但多寫了兩
行：一行負責在迴圈開始前，先把\texttt{window}設成 0；一行負責在迴圈裡手動把
\texttt{window}加 1。多這兩行不只是麻煩，還多了一個風險：如果哪天不小心漏寫了
\texttt{window += 1}這一行，這段程式碼就會變成無窮迴圈，跟第~\ref{sec:while-basics}~
節看過的坑一模一樣。同樣的重複，用\texttt{for}搭配\texttt{range()}既不用自己管理計數
變數，也不用擔心忘記更新它。

\keypoint{已知次數、或是走訪一批固定內容的重複，優先用\texttt{for}：\texttt{range()}
和直接走訪序列都會替你處理好「這一輪走到哪、還有沒有下一輪」，不會有忘記更新變數導致
無窮迴圈的風險。只有在停止條件真的要看執行過程中才會變化的狀態（使用者輸入、還沒存夠
的金額）時，才需要\texttt{while}。}

這也是為什麼你會發現，越到後面的章節，Python 程式碼裡\texttt{for}出現的頻率遠高於
\texttt{while}。第 \ref{ch:lists} 章要學的 list、第 \ref{ch:dict-set} 章的 dict，只要是「有一批東西，要對每一個都
做一件事」的情境，幾乎都是先想\texttt{for}，很少會拿\texttt{while}硬幹。

原因跟這一節
比較\texttt{windows.py}／\texttt{windows-while.py}時看到的一樣：\texttt{for}走訪容器
不用你自己管索引、不用自己設計停止條件，程式碼更短，也更不容易寫出無窮迴圈。
\texttt{while}真正無可取代的地方，是那些連你自己都不知道要跑幾輪的情境，像使用者輸入、
或是之後會遇到的網路請求重試這一類。

\begin{drills}[label=drills:for-vs-while]
\item 下面這個任務比較適合用\texttt{for}還是\texttt{while}？「檢查購物清單上的每一項，
在清單上打勾」（提示：你在開始檢查之前，知不知道清單上有幾項？）
\item 下面這個任務比較適合用\texttt{for}還是\texttt{while}？「打電話給客服，一直等到
有人接聽為止」
\item 抓蟲：把\texttt{windows-while.py}裡的\texttt{window += 1}拿掉，會發生什麼事？
\item 第~\ref{sec:while-basics}~節的\texttt{saving.py}用\texttt{while}存錢存到存夠為
止。如果需求換成「固定存 30 天，不管有沒有存到目標」，這時候比較適合用\texttt{for}還是
\texttt{while}？為什麼？
\end{drills}

\section{提早結束：\texttt{break} 與線性搜尋}
\label{sec:break-search}

前面幾節的迴圈，都是乖乖跑完該跑的次數，或是跑到條件不成立才停。但有些情境不一樣。你在
翻找家裡的抽屜找鑰匙時，只要找到了，就不用繼續翻剩下的抽屜；反過來，如果翻遍了所有抽屜
都沒找到，你才會下結論說「鑰匙不在這裡」。這種「找到就提早結束，找完都沒找到才下結論」
的動作，Python 用\texttt{break}處理。
\index{break}

\begin{code}{python}{for 搭配 break 提早結束}
for item in sequence:
    if condition:
        break
    repeated_action
\end{code}

一旦執行到\texttt{break}，Python 會立刻跳出目前這個迴圈，後面還沒跑的輪次全部跳過，
直接接著執行迴圈後面的程式碼。

\begin{example}{python}{(find-key.py) 翻找抽屜找鑰匙}
key_location = 3
found = False

for drawer in range(1, 6):
    print(f"Searching drawer {drawer}")
    if drawer == key_location:
        found = True
        break

if found:
    print(f"Key found in drawer {key_location}!")
else:
    print("Searched every drawer, no key found")
\end{example}

\iosample{}{Searching drawer 1 \\ Searching drawer 2 \\ Searching drawer 3 \\ Key found in
drawer 3! \\}

這種「從頭開始一個一個找，找到就停，找完都沒找到才說沒有」的做法，有個正式的名字叫
\textbf{線性搜尋}。
\index{線性搜尋}
上面的例子用了一個\texttt{found}變數記錄「到底有沒有找到」，等迴圈
結束（不管是提早\texttt{break}結束，還是正常跑完）才根據\texttt{found}的值決定要印哪一
句話。這個順序很重要，接下來要示範一個新手很常見的寫法錯誤。

\begin{example}{python}{(find-key-bug.py) 把兩種結果都寫在迴圈裡面判斷}
key_location = 5

for drawer in range(1, 6):
    if drawer == key_location:
        print("Key found!")
    else:
        print("No key in this drawer, search ended")
    break
\end{example}

\iosample{}{No key in this drawer, search ended \\}

鑰匙明明放在第 5 個抽屜，這段程式碼卻在翻第 1 個抽屜就說「search ended」，而且完全沒提到
鑰匙到底找到了沒有。問題出在\texttt{break}的縮排位置：它跟\texttt{if}／\texttt{else}
同一層，不管這一輪翻到的是不是目標抽屜，迴圈都只跑一輪就會被\texttt{break}停掉。

正確的
寫法是像上面\texttt{find-key.py}那樣，只有真的符合條件（\texttt{drawer == key\_location}）
的那一輪，才把\texttt{break}放進\texttt{if}裡面。沒找到的那幾輪只該正常結束這一輪、
繼續翻下一個抽屜，不該提前喊停，也不該在迴圈裡面就直接下「沒找到」的結論。

\keypoint{\texttt{break}只會讓你跳出目前正在跑的這一層迴圈，不會讓整支程式停下來。
程式會接著執行迴圈後面的程式碼。}

順帶一提，Python 內建的\texttt{in}運算子，其實就是幫你把整個「一個一個找、找到就算」
的搜尋迴圈包好了：\texttt{5 in range(1, 6)}會直接告訴你 True，不用自己手寫迴圈。之後
遇到單純只想知道「有沒有」而不需要知道「翻了幾次」的情境，可以直接用\texttt{in}；但
搜尋迴圈本身的骨架，還是值得你自己動手寫過一次，之後才看得懂\texttt{in}幫你省了什麼。
\index{in 運算子}

\begin{drills}[label=drills:break-search]
\item 不要真的執行，用眼睛看：下面這段程式碼會印出哪幾行？
\begin{verbatim}
target = 4

for number in range(1, 7):
    if number == target:
        print(f"Found {number}")
        break
    print(f"{number} is not the target")
\end{verbatim}
\item 抓蟲：下面這段程式碼哪裡有問題？為什麼鑰匙明明放在第 2 個抽屜，程式卻說沒找到？
\begin{verbatim}
key_location = 2

for drawer in range(1, 4):
    if drawer == key_location:
        print("Key found!")
    else:
        print("No key in this drawer, search ended")
    break
\end{verbatim}
\item 如果找鑰匙的情境改成巢狀的（例如每個房間裡都有好幾個抽屜，要先走訪房間、再走訪
抽屜），只在最內層（抽屜那層）寫\texttt{break}，找到鑰匙後真的能讓兩層迴圈都停下來嗎？
\end{drills}

\section{跳過本輪：\texttt{continue}}
\label{sec:continue}

\texttt{break}是「整個迴圈都不要繼續了」，但有時候你只是想跳過某一輪，其他輪照常進行。
檢查一週七天的運動紀錄時，週日是休息日，遇到週日就跳過不用記錄，但接下來的每一週還是要
繼續檢查。這種「跳過本輪、直接進到下一輪」的動作，Python 用\texttt{continue}處理：一旦
執行到\texttt{continue}，這一輪縮排在它後面的程式碼全部不執行，直接跳到下一輪重新開始。
\index{continue}

\begin{example}{python}{(weekly-check.py) 檢查一週七天的運動紀錄}
for day in range(1, 8):
    if day == 7:
        continue
    print(f"Day {day}: workout complete")

print("Weekly check complete")
\end{example}

\iosample{}{Day 1: workout complete \\ Day 2: workout complete \\ Day 3: workout complete
\\ Day 4: workout complete \\ Day 5: workout complete \\ Day 6: workout complete
\\ Weekly check complete \\}

留意第 7 天完全沒有出現在輸出裡。不是程式漏掉了，是\texttt{continue}讓那一輪直接跳過
了底下的\texttt{print}，但迴圈本身並沒有結束，第 7 天結束後迴圈也正常跑完，接著印出
「Weekly check complete」。

\begin{drills}[label=drills:continue]
\item 不要真的執行，用眼睛看：\texttt{for n in range(1, 6):}底下接
\texttt{if n == 3: continue}，再接\texttt{print(n)}，會印出哪幾個數字？
\item 如果把上面\texttt{weekly-check.py}裡判斷的條件從\texttt{day == 7}改成
\texttt{day == 3}，\texttt{continue}換成\texttt{break}之後，輸出會有什麼不一樣？
\item 抓蟲（延伸）：如果一個巢狀迴圈的內層寫了\texttt{continue}，外層那一輪會不會也
被跳過？
\end{drills}

\section{迴圈正常跑完才做的事：for-else 與 while-else}
\label{sec:loop-else}

第~\ref{sec:break-search}~節翻抽屜找鑰匙的例子，用了一個\texttt{found}變數記錄「到底
有沒有找到」，等迴圈結束才根據\texttt{found}的值決定要印哪一句話。Python 其實提供一個
內建的語法，可以省掉這個額外的變數：\texttt{for}和\texttt{while}都可以接一個\texttt{else}
區塊，這個\texttt{else}只有在迴圈正常跑完、沒有被\texttt{break}提早結束的情況下才會
執行。
\index{for-else／while-else}

\begin{code}{python}{for-else：迴圈正常跑完才執行 else}
for item in sequence:
    if condition:
        break
else:
    action_if_not_broken
\end{code}

留意\texttt{else}的縮排跟\texttt{for}同一層，不是跟裡面的\texttt{if}同一層。這裡的
\texttt{else}問的其實是「整個迴圈跑完的過程中，有沒有被\texttt{break}打斷過」，跟裡面
那個\texttt{if}有沒有成立是兩回事。只要迴圈從頭到尾都沒碰到\texttt{break}，不管是正常
跑滿所有輪次，還是\texttt{sequence}根本是空的、一次都沒跑，\texttt{else}都會執行；只要
中途執行到\texttt{break}，\texttt{else}就會被跳過。

\begin{example}{python}{(find-key-else.py) 用 for-else 改寫翻抽屜找鑰匙}
key_location = 3

for drawer in range(1, 6):
    print(f"Searching drawer {drawer}")
    if drawer == key_location:
        print(f"Key found in drawer {drawer}!")
        break
else:
    print("Searched every drawer, no key found")
\end{example}

\iosample{}{Searching drawer 1 \\ Searching drawer 2 \\ Searching drawer 3 \\ Key found in
drawer 3! \\}

跟第~\ref{sec:break-search}~節的\texttt{find-key.py}比較一下：這裡完全不需要
\texttt{found}這個變數，迴圈找到鑰匙就\texttt{break}、印出訊息；找不到的話，迴圈正常
跑完六輪，\texttt{else}就會接手印出「Searched every drawer, no key found」。少了一個額外變數，
邏輯卻一樣清楚。

\begin{example}{python}{(find-target-else.py) 迴圈正常跑完，else 被執行}
for number in range(1, 4):
    if number == 5:
        print("Found 5")
        break
else:
    print("5 not found")
\end{example}

\iosample{}{5 not found \\}

\texttt{range(1, 4)}裡從頭到尾都不會出現 5，迴圈跑完三輪，一次\texttt{break}都沒碰到，
於是\texttt{else}被執行，印出「5 not found」。

\keypoint{很多人第一次看到\texttt{for-else}會誤以為\texttt{else}是「迴圈一次都沒跑」
才執行，其實不是。\texttt{else}只跟「有沒有被\texttt{break}打斷」有關，跟迴圈本身跑了
幾輪、甚至有沒有跑完全沒關係，只要正常跑完（哪怕跑了零輪），\texttt{else}就會執行。}

\texttt{while}也可以接同樣的\texttt{else}，規則完全一樣：\texttt{while}的條件變成不
成立、正常結束迴圈，\texttt{else}就會執行；中途被\texttt{break}提早結束，\texttt{else}
就會被跳過。

\begin{code}{python}{while-else：迴圈正常結束才執行 else}
while condition:
    if early_stop_condition:
        break
else:
    action_if_not_broken
\end{code}

\begin{example}{python}{(login-attempts.py) 密碼嘗試三次後鎖定帳號}
attempts = 0
password = "1234"
guess = "0000"

while attempts < 3:
    if guess == password:
        print("Password correct")
        break
    print("Wrong password, try again")
    attempts += 1
else:
    print("Max attempts reached, account locked")
\end{example}

\iosample{}{Wrong password, try again \\ Wrong password, try again \\ Wrong password, try again
\\ Max attempts reached, account locked \\}

這裡\texttt{guess}是先寫死的示範值，只是要展示整個流程；如果要讓使用者真的每一次都
重新輸入，可以把\texttt{guess = input(...)}搬進迴圈裡面，就像第~\ref{sec:while-input}~
節做過的那樣。迴圈跑了三輪，\texttt{guess}始終不等於\texttt{password}，
\texttt{attempts}從 0 累加到 3，\texttt{while attempts < 3}的條件變成不成立，迴圈正常
結束，不是被\texttt{break}打斷，於是\texttt{else}執行，印出鎖定訊息。

\stylepoint{\texttt{for-else}／\texttt{while-else}是 Python 一個相對少見的語法，很多
其他程式語言完全沒有這個概念，所以就算你之後查資料查到別的語言的範例，也不用太意外看
不到它。適合用在「找到就\texttt{break}，沒找到才要做某件事」這種場景，能省掉一個額外
的旗標變數；但如果你或合作的人對這個語法不熟，用一個明確的旗標變數（像
\texttt{find-key.py}那樣）有時候反而更好讀，兩種寫法都要看得懂。}

\begin{drills}[label=drills:loop-else]
\item 不要真的執行，用眼睛看：下面這段程式碼會印出哪些行？
\begin{verbatim}
for number in range(1, 5):
    if number == 3:
        print("Found 3")
        break
    print(f"{number} is not the target")
else:
    print("Not found")
\end{verbatim}
\item 如果把上面的\texttt{number == 3}改成\texttt{number == 9}，輸出會有什麼不同？
\item 抓蟲：下面這段\texttt{while-else}，想在猜對密碼時印出「Password correct」，猜錯三次印出
「Account locked」，但實際執行會發現兩句話都印出來了，問題出在哪？
\begin{verbatim}
attempts = 0
while attempts < 3:
    if guess == password:
        print("Password correct")
    attempts += 1
else:
    print("Account locked")
\end{verbatim}
（提示：對照\texttt{login-attempts.py}，這裡少了什麼關鍵字？）
\end{drills}

\section{三種常見的迴圈任務長相：累加、計數、搜尋}
\label{sec:accumulate-count}

回頭看看這一章寫過的幾個例子：

\begin{itemize}
\item 第~\ref{sec:while-basics}~節的存錢範例，是把每天存的錢一次一次加總起來，直到存
夠為止。
\item 第~\ref{sec:for-string}~節數空白字元的範例，是每遇到一次符合條件的字元，就把計
數往上加一。
\item 第~\ref{sec:break-search}~節翻抽屜找鑰匙，是一個一個找，找到就停。
\end{itemize}

這三個例子表面上做的事情完全不一樣，但拆開來看，骨架其實長得很像：

\begin{code}{python}{迴圈裡的常見任務骨架}
result = initial_value    # 迴圈開始前，先準備好需要的變數

for item in sequence:
    if condition:
        result = new_result_value   # 迴圈中，視情況每一輪更新它

print(result)              # 迴圈結束後，檢查最終結果
\end{code}

三種長相的差別，只在「迴圈中每一輪具體做了什麼」：

\begin{itemize}
\item \textbf{累加}：\index{累加}每一輪把值加到\texttt{result}上，像第~\ref{sec:while-basics}~節
的存錢範例，\texttt{savings += daily\_amount}。
\item \textbf{計數}：\index{計數}每一輪遇到符合條件的情況，就把\texttt{result}加 1，像
第~\ref{sec:for-string}~節數空白字元，\texttt{space\_count += 1}。
\item \textbf{搜尋}：\index{搜尋（迴圈任務長相）}每一輪檢查是不是找到目標，找到就用\texttt{break}提早結束，像
第~\ref{sec:break-search}~節翻抽屜找鑰匙。
\end{itemize}

再用一個新例子練一次計數模式：

\begin{example}{python}{(count-vowels.py) 數一句話裡有幾個母音字母}
sentence = "today is a good day"
vowels = "aeiou"
vowel_count = 0

for char in sentence:
    if char in vowels:
        vowel_count += 1

print(f"This sentence has {vowel_count} vowels")
\end{example}

\iosample{}{This sentence has 7 vowels \\}

這三種骨架都是「先準備一個變數，每一輪視情況更新它」，但\texttt{result}不一定要是
數字。同一套邏輯換成字串一樣能用：先準備一個空字串，每一輪把新的內容用串接的方式接
上去，就能一路把過程中產生的東西收集起來。

\begin{example}{python}{(mask-spaces.py) 把一句話裡的空白換成底線，重新組成新字串}
sentence = "today is a good day"
report = ""

for char in sentence:
    if char == " ":
        report += "_"
    else:
        report += char

print(report)
\end{example}

\iosample{}{today\_is\_a\_good\_day \\}

\texttt{report}一開始是空字串，每一輪迴圈都用\texttt{+=}把一個字元（或代替空白的底線）
接到\texttt{report}後面，跑完整個\texttt{sentence}，\texttt{report}就變成一句話從頭到
尾都處理過一遍的新字串。這正是累加骨架的另一種長相，只是\texttt{result}累加的不是數字，
是字串。

之後學到 list，你會看到同一個「先準備一個空的容器，每一輪塞一點東西進去」的
邏輯，只是換成 list 搭配\texttt{append()}方法。不管容器是數字、字串，還是之後的 list，
先準備、逐輪更新、最後看結果，這個節奏都是一樣的。

\keypoint{累加、計數、搜尋這三種骨架，會在你之後幾乎每一章、每一支程式裡反覆出現。
「迴圈前先準備好一個變數，迴圈中每一輪視情況更新它，迴圈後看最終結果」這個順序，值得你
牢牢記住，不管\texttt{result}放的是數字還是字串，節奏都一樣。之後遇到一個新問題時，先
問自己一句：「這比較像累加、計數、還是搜尋？」通常就能很快想出迴圈該怎麼寫。}

\begin{drills}[label=drills:accumulate-count]
\item 不要真的執行，用眼睛看：下面這段程式碼最後印出的\texttt{total}是多少？
\begin{verbatim}
total = 0
for i in range(1, 5):
    total += i
print(total)
\end{verbatim}
\item 填空：讓下面這段程式碼算出\texttt{word}這個字串裡有幾個字母\texttt{"e"}。
\begin{verbatim}
word = "elephant"
count = ____

for letter in word:
    if letter == "e":
        count ____

print(count)
\end{verbatim}
\item 填空：想要用累加字串的方式，把\texttt{word = "python"}裡的每個字母都重複兩次接
起來（結果應該是\texttt{"ppyytthhoonn"}），迴圈該怎麼寫？（提示：可以用兩次\texttt{+}
把同一個字母接兩遍。）
\end{drills}

\section*{練習解答}
\addcontentsline{toc}{section}{練習解答}

\paragraph{第~\ref{sec:for-range}~節}
\begin{enumerate}
\item 印出：0、1、2（三行）。
\item 8，\texttt{range(2, 9, 2)}依序是 2、4、6、8，9 本身不會被包含進去。
\item \texttt{range(9, 0, -2)}，從 9 開始，每次減 2，數到比 1 大的地方為止，最後一個
數字剛好是 1。
\item 0 次。\texttt{range(5, 5)}起點跟終點相同，代表沒有任何數字落在這個範圍內，迴圈
本體一次都不會執行。
\end{enumerate}

\paragraph{第~\ref{sec:for-string}~節}
\begin{enumerate}
\item 印出：C、A、T（三行），\texttt{.upper()}把每個字元轉成大寫。
\item 缺少冒號，錯誤訊息會是\texttt{SyntaxError: expected ':'}。修法是在
\texttt{"abc"}後面補上冒號：\texttt{for char in "abc":}。
\end{enumerate}

\paragraph{第~\ref{sec:nested-loops}~節}
\begin{enumerate}
\item 6 行。外層迴圈跑 2 輪，內層迴圈每一輪都完整跑 3 輪，2 乘 3 等於 6。
\item 10 次。外層\texttt{range(1, 6)}是 5 個房間，內層\texttt{range(1, 3)}是每個房間
2 扇窗戶，5 乘 2 等於 10。
\item 只在內層寫\texttt{break}不夠，因為\texttt{break}只會跳出離它最近的那一層迴圈
（抽屜／窗戶那一層）。想讓外層也跟著停，需要額外設一個變數（例如\texttt{found}）記錄
「已經找到了」，在外層迴圈裡也檢查這個變數，符合條件時外層自己也執行一次
\texttt{break}。
\end{enumerate}

\paragraph{第~\ref{sec:enumerate}~節}
\begin{enumerate}
\item 印出：\texttt{0: c}、\texttt{1: a}、\texttt{2: t}（三行）。
\item \texttt{for index, char in enumerate(word): print(f"Letter \{index + 1\} is
\{char\}")}，跟\texttt{range()}一樣，索引從 0 開始算，要顯示「第幾個」給人看時，記得
補上\texttt{+ 1}。
\item 少了\texttt{break}。找到\texttt{"p"}就印出索引之後沒有提早結束迴圈，所以
\texttt{"apple"}裡兩個\texttt{"p"}（索引 1 跟 2）都會被印出來，不只印出第一次出現的
位置。
\end{enumerate}

\paragraph{第~\ref{sec:while-basics}~節}
\begin{enumerate}
\item 印出 3 次「chances left」（3 chances left、2 chances left、1 chances left），因為
\texttt{n}每輪都有確實減 1，跑到\texttt{n == 0}時條件就不成立了。
\item 這是無窮迴圈，迴圈裡從來沒有更新\texttt{attempts}，條件永遠成立。修法是在迴圈裡
補上一行讓\texttt{attempts}遞減，例如\texttt{attempts -= 1}。
\item 這也是無窮迴圈：條件用的是\texttt{remaining}，但迴圈裡實際更新的卻是
\texttt{count}，\texttt{remaining}從頭到尾都沒變過，\texttt{remaining > 0}永遠成立。
修法是把\texttt{count += 1}改成\texttt{remaining -= 1}，或是額外補上更新
\texttt{remaining}的那一行。
\item 不會是無窮迴圈。6 每次減 2 會剛好經過 0（6、4、2、0），條件\texttt{n != 0}在
\texttt{n}變成 0 的那一輪開始前就不成立了，迴圈正常停止，印出 6、4、2 三行。
\end{enumerate}

\paragraph{第~\ref{sec:while-input}~節}
\begin{enumerate}
\item 1 行，「Logged: mop floor」。「mop floor」不是結束關鍵字，所以迴圈印出這一行、再問一次；
下一次輸入「done」，條件不成立，迴圈結束，不會再印出「Logged:」開頭的行。
\item 不會，因為第一次讀到的輸入就已經是「done」，\texttt{while}的條件一開始就不成立，
迴圈本體一次都不會執行，只會印出「Chore log complete!」。
\end{enumerate}

\paragraph{第~\ref{sec:for-vs-while}~節}
\begin{enumerate}
\item 適合用\texttt{for}，因為開始檢查之前，購物清單上有幾項是已知的、固定的。
\item 適合用\texttt{while}，因為你不知道要等幾秒、等幾次才會有人接聽，只知道「有人接
才停」這個條件。
\item 條件永遠成立，\texttt{window < 5}變成無窮迴圈，程式會不斷印出「Cleaning window
1」，因為\texttt{window}從頭到尾都是 0，永遠小於 5。
\item 適合用\texttt{for}。「固定存 30 天」在動手之前就已經確定次數，屬於已知範圍，改用
\texttt{for i in range(30):}最直接；換成\texttt{while}反而得自己額外維護一個計數變數，
多繞一圈。
\end{enumerate}

\paragraph{第~\ref{sec:break-search}~節}
\begin{enumerate}
\item 印出：1 is not the target、2 is not the target、3 is not the target、Found 4，迴圈
跑到\texttt{number == 4}時符合條件，印出訊息後用\texttt{break}提早結束，後面的 5、6
都不會再檢查。
\item \texttt{break}的縮排跟\texttt{if}／\texttt{else}同一層，不管有沒有符合條件，
迴圈都只跑一輪就會被強制停掉，翻第 1 個抽屜沒找到，就直接被\texttt{break}停住，根本
沒機會翻到第 2 個抽屜。修法是把\texttt{break}移進\texttt{if}裡面，只有真的找到才提早
結束。
\item 不能。只在最內層（抽屜層）寫\texttt{break}，只會讓那一個房間的抽屜迴圈提早結束，
外層的房間迴圈完全不受影響，還是會繼續進到下一個房間重新翻抽屜。想讓兩層都停，得額外
處理外層（參考第~\ref{sec:nested-loops}~節的做法）。
\end{enumerate}

\paragraph{第~\ref{sec:continue}~節}
\begin{enumerate}
\item 印出：1、2、4、5，\texttt{n == 3}那一輪被\texttt{continue}跳過，沒有執行到
\texttt{print(n)}。
\item 用\texttt{continue}的話，只有第 3 天被跳過，第 4 到第 7 天照常印出，輸出還是有
六行「workout complete」；換成\texttt{break}的話，迴圈跑到第 3 天就整個結束，第 4 到第 7 天
根本不會再檢查，只會印出「Day 1」「Day 2」兩行。\texttt{continue}是「跳過這一輪，
其他輪照常進行」，\texttt{break}是「後面全部都不要了」，兩者不一樣。「Weekly check
complete」那一行則不管哪種寫法都會印出來，因為它在迴圈外面，迴圈結束（不管是正常跑完還是被
\texttt{break}提早結束）之後都會執行到它。
\item 不會。\texttt{continue}只會跳過內層迴圈自己這一輪剩下的動作，外層迴圈完全不知道
內層發生了什麼事，會繼續照常進行自己的輪次。
\end{enumerate}

\paragraph{第~\ref{sec:loop-else}~節}
\begin{enumerate}
\item 印出：1 is not the target、2 is not the target、Found 3。因為找到 3 就
\texttt{break}，\texttt{else}被跳過，不會印出「Not found」。
\item 迴圈會跑完 1 到 4 全部四輪，一次都沒有符合\texttt{number == 9}，也就是一次
\texttt{break}都沒碰到，於是\texttt{else}執行，最後多印出「Not found」這一行，其餘四行
「is not the target」照常印出。
\item 少了\texttt{break}。猜對密碼只有印出「Password correct」，但迴圈沒有提早結束，還是會把
剩下的輪次跑完，跑完之後（沒有被\texttt{break}打斷）\texttt{else}一定會執行，所以
「Account locked」永遠都會印出來，不管密碼有沒有猜對。
\end{enumerate}

\paragraph{第~\ref{sec:accumulate-count}~節}
\begin{enumerate}
\item 10，\texttt{range(1, 5)}依序是 1、2、3、4，加總起來是 1+2+3+4=10。
\item \texttt{count = 0}，以及\texttt{count += 1}，先把\texttt{count}初始化成 0，
每遇到一個字母\texttt{"e"}就加 1。
\item \texttt{result = ""}，接著\texttt{for letter in word: result += letter +
letter}，每一輪把同一個字母接兩次，串到\texttt{result}後面。
\end{enumerate}
